\documentclass{beamer}
\usetheme{Madrid}
\usecolortheme{dolphin}
\definecolor{charcoal}{RGB}{54,69,79}
\setbeamercolor{structure}{fg=charcoal}
\setbeamercolor{palette primary}{bg=charcoal,fg=white}
\setbeamercolor{palette secondary}{bg=charcoal!90,fg=white}
\setbeamercolor{palette tertiary}{bg=charcoal!80,fg=white}
\setbeamercolor{palette quaternary}{bg=charcoal!70,fg=white}
\setbeamercolor{frametitle}{bg=charcoal,fg=white}
\setbeamercolor{title}{bg=charcoal,fg=white}
\usepackage{hyperref}
\usepackage{graphicx}
\usepackage{booktabs}
\usepackage{listings}
\usepackage{color}
\definecolor{dkgreen}{rgb}{0,0.6,0}
\definecolor{gray}{rgb}{0.5,0.5,0.5}
\definecolor{mauve}{rgb}{0.58,0,0.82}
\lstset{
  language=Java,
  aboveskip=3mm,
  belowskip=3mm,
  showstringspaces=false,
  columns=flexible,
  basicstyle={\small\ttfamily},
  numbers=none,
  keywordstyle=\color{blue},
  commentstyle=\color{dkgreen},
  stringstyle=\color{mauve},
  breaklines=true,
  breakatwhitespace=true,
  tabsize=2
}
\title[Abstract Classes \& Interfaces]{Abstract Classes and Interfaces}
\subtitle{COP2250: Java Programming}
\author{Kevin Pyatt, Ph.D.}
\institute{State College of Florida \\ Pyatt Labs}
\date{Week 14}
\begin{document}
\begin{frame}
  \titlepage
\end{frame}
\begin{frame}{Today's Objectives}
\begin{itemize}
  \item Declare and use abstract classes
  \item Write abstract methods subclasses must implement
  \item Define and implement interfaces
  \item Understand when to use abstract class vs interface
  \item Apply both to a shape hierarchy
\end{itemize}
\end{frame}
\begin{frame}[fragile]{Abstract Class}
\textbf{Cannot be instantiated. Subclasses must implement abstract methods.}
\begin{lstlisting}
abstract class GeometricObject {
    private String color;

    public GeometricObject(String color) {
        this.color = color;
    }

    public abstract double getArea();
    public abstract double getPerimeter();

    public String toString() {
        return "Color: " + color
             + "\nArea: " + getArea();
    }
}
\end{lstlisting}
\textit{toString() calls getArea() --- the subclass version runs. Polymorphism.}
\end{frame}
\begin{frame}[fragile]{Interface}
\textbf{A pure contract. Implementing class must provide all methods.}
\begin{lstlisting}
interface Colorable {
    void howToColor();
}
\end{lstlisting}
\vspace{0.3cm}
\begin{tabular}{ll}
\toprule
\textbf{Abstract Class} & \textbf{Interface} \\
\midrule
Can have fields & No instance fields \\
Can have constructors & No constructors \\
Can have concrete methods & Methods abstract by default \\
Extend one only & Implement many \\
\bottomrule
\end{tabular}
\end{frame}
\begin{frame}[fragile]{Implementing Both}
\begin{lstlisting}
class Circle extends GeometricObject implements Colorable {
    private double radius;

    public Circle(double radius, String color) {
        super(color);
        this.radius = radius;
    }

    public double getArea() { return Math.PI * radius * radius; }
    public double getPerimeter() { return 2 * Math.PI * radius; }

    public void howToColor() {
        System.out.println("Color the circle with radius "
            + radius);
    }
}
\end{lstlisting}
\textit{Must implement all abstract methods from GeometricObject and Colorable.}
\end{frame}
\begin{frame}[fragile]{Polymorphism and Casting}
\begin{lstlisting}
GeometricObject g = new Circle(5.0, "red");
System.out.println(g.getArea());   // Circle's getArea()

// Cast required to call interface method
((Colorable) g).howToColor();
\end{lstlisting}
\vspace{0.3cm}
\begin{itemize}
  \item Reference type: \texttt{GeometricObject}
  \item Actual object: \texttt{Circle}
  \item Cast needed only to access \texttt{Colorable} methods
  \item Wrong cast at runtime: \texttt{ClassCastException}
\end{itemize}
\end{frame}
\begin{frame}{Abstract Class vs Interface}
\textbf{Use abstract class when:}
\begin{itemize}
  \item Subclasses share fields or concrete methods
  \item Strong is-a relationship
\end{itemize}
\vspace{0.3cm}
\textbf{Use interface when:}
\begin{itemize}
  \item Unrelated classes share behavior
  \item Multiple contracts needed
\end{itemize}
\vspace{0.3cm}
\textit{Colorable applies to Circle, Rectangle, even Car. They share behavior, not ancestry.}
\end{frame}
\begin{frame}{Key Terms}
\begin{description}
  \item[Abstract class] Cannot be instantiated; may contain abstract methods
  \item[Abstract method] No body; subclass must implement
  \item[Interface] Pure contract; defines method signatures only
  \item[\texttt{implements}] Keyword a class uses to fulfill an interface
  \item[Cast] Treating a reference as a specific type
  \item[\texttt{instanceof}] Tests whether an object is of a specific type
\end{description}
\end{frame}
\begin{frame}{Common Mistakes}
\begin{tabular}{ll}
\toprule
\textbf{Mistake} & \textbf{Fix} \\
\midrule
Instantiating abstract class & Use a concrete subclass \\
Missing \texttt{super(color)} & Parent constructor must be called \\
Interface method not implemented & Compiler tells you which one \\
Cast exception at runtime & Check with \texttt{instanceof} first \\
Abstract method has a body & Remove body, add semicolon \\
\bottomrule
\end{tabular}
\end{frame}
\begin{frame}{Lab --- Abstract Shape Hierarchy}
\textbf{Build in order:}
\begin{enumerate}
  \item Define \texttt{Resizable} interface: \texttt{void resize(double factor)}
  \item Abstract \texttt{Shape}: color, abstract \texttt{getArea()}, abstract \texttt{getPerimeter()}
  \item \texttt{Circle extends Shape implements Resizable}
  \item \texttt{Rectangle extends Shape implements Resizable}
  \item Test: create both, print area, resize, print again
\end{enumerate}
\vspace{0.3cm}
\textbf{Demo Requirement:} Compile live, explain abstract vs interface.\\
\vspace{0.3cm}
\textbf{Next:} Assignment 13 --- Abstract Classes and Interfaces
\end{frame}
\end{document}
